\section{Schedule management}
\label{sec:schedule}

The campaign schedule turns the WBS into a dated baseline. Dates come from a finish-to-start activity-on-node network on a seven-day pad calendar, 10~August to 1~December~2026, because LOX load, static fire and lift-off do not stop for weekends. Baseline $T$-0 is 15~November~2026, inside the charter's Q1~2027 window. A five-day pull-in to 10~November exists only if R-14 is realised. It is not the baseline.

\subsection{Network logic and float analysis}
\label{sec:schedule-network}

Links are finish-to-start with zero lag, drawn activity-on-node. Textbook discrete-time notation writes $EF = ES + D$ when $EF$ is exclusive. Every date on the Bowen calendar is inclusive. The identities used here are
\begin{equation}
ES_{i} = \max_{j \in P_{i}}(EF_{j}) + 1,
\qquad
EF_{i} = ES_{i} + D_{i} - 1,
\label{eq:fwd}
\end{equation}
\begin{equation}
LF_{i} = \min_{k \in S_{i}}(LS_{k}) - 1,
\qquad
LS_{i} = LF_{i} - D_{i} + 1,
\label{eq:bwd}
\end{equation}
$P_{i}$ and $S_{i}$ are the predecessor and successor sets of activity $i$. Where $P_{i}$ is empty, $ES$ is the campaign start. Total float and free float are
\begin{equation}
TF_{i} = LS_{i} - ES_{i} = LF_{i} - EF_{i},
\qquad
FF_{i} = \min_{k \in S_{i}}(ES_{k}) - EF_{i} - 1,
\label{eq:float}
\end{equation}
If $S_{i}$ is empty, $FF_{i} = TF_{i}$. Launch-critical means $TF_{i} = 0$ on a path that sets $T$-0 or closeout. Table~\ref{tab:precedence} is the full pass. It matches the locked register.

The forward pass gives two $TF = 0$ chains. They stay independent until hot fire.
\begin{itemize}[leftmargin=1.4em,itemsep=0.25em]
\item Vehicle path: A-100 $\rightarrow$ A-122 $\rightarrow$ A-123 $\rightarrow$ A-132 $\rightarrow$ A-133. A-100 is the external \$0 freight wait, not a Gilmour launch-critical bar.
\item Permit path: ASA / CASA / GBRMPA permits (A-113) into static fire (A-133).
\end{itemize}
RF calibration (A-125) feeds static fire, so the array is live before hot fire. It is not launch-critical. After static fire, A-134 through closeout (A-145) is one critical chain. The PEP lock (A-111) finishes on 2~October by an imposed Must-Finish-On (Gate~1). It is not launch-critical. The Week~4 charter approved the document, not field mobilisation.

The paths meet in two staged merges.
\begin{itemize}[leftmargin=1.4em,itemsep=0.25em]
\item Wet Dress Rehearsal waits on stacking, cryogenic GSE (A-124), avionics power-on (A-131) and Juru monitors (A-114). That is the technical merge.
\item Static fire then waits on the rehearsal, flight permits \emph{and} RF calibration. That is the regulatory and range-instrument merge.
\end{itemize}
Launch Readiness Review (LRR, A-141) waits on post-fire diagnostics, not on static fire. Signing that flight-clearance gate on incomplete hot-fire data would break the hold-point logic of an AS9100D system \citep{sae2016}.

Feeder float is not interchangeable.
\begin{itemize}[leftmargin=1.4em,itemsep=0.25em]
\item Avionics has one day into the rehearsal.
\item Cryogenic GSE has three days.
\item RF calibration has 39 days. It sits in August--September so it does not sit on the October pad peak. Moving it does not change 15~November.
\item Cleanroom mobilisation (A-121) has $TF = 3$ but $FF = 0$, so delay is absorbed before $T$-0 only by consuming the GSE float.
\end{itemize}
Total float protects $T$-0. Free float protects the successor. Figure~\ref{fig:network_logic} draws the links.

\begin{figure}[H]
\centering
\includegraphics[width=\linewidth]{network_logic.pdf}
\caption[Activity-on-node network]{Activity-on-node network for the Bowen campaign. Node bands are $ES$ / ID / $EF$, $LS$ / duration / $LF$, and $TF$. Crimson arrows are launch-critical finish-to-start links. The vehicle path and the permit path merge at static fire; RF calibration is the labelled south-routed feeder into the same merge; LRR waits on post-fire diagnostics, not on static fire. The PEP lock is Gate-1 constrained and is not launch-critical. Freight receipt is a grey dashed external \$0 driver, not a crimson Gilmour bar.}
\label{fig:network_logic}
\end{figure}

\begin{landscape}
\small
\setlength{\tabcolsep}{4.2pt}
\renewcommand{\arraystretch}{1.15}
\begin{longtable}{@{}l l l r l c c c c r r c@{}}
\caption[Activity precedence and float register]{Master campaign activity precedence and float register. Dates are inclusive. $FF = \min(ES_{\text{succ}}) - EF - 1$ days, or $FF = TF$ if there is no successor. The PEP lock $TF = 0$ is an imposed Gate~1 finish, not launch-critical. The freight wait $TF = 0$ is an external pin (CP $=$ Ext), not a Gilmour launch-critical bar.}
\label{tab:precedence}\\
\toprule
ID & WBS & Activity & $D$ & Pred. & ES & EF & LS & LF & TF & FF & CP \\
\midrule
\endfirsthead
\caption[]{Master campaign activity precedence and float register (continued).}\\
\toprule
ID & WBS & Activity & $D$ & Pred. & ES & EF & LS & LF & TF & FF & CP \\
\midrule
\endhead
\bottomrule
\endfoot
A-100 & EXT & Vehicle receipt (external driver) & 35 & --- & 10 Aug & 13 Sep & 10 Aug & 13 Sep & 0 & 0 & Ext \\
A-111 & 1.1.1 & PEP and integration baseline & 54 & --- & 10 Aug & 2 Oct & 10 Aug & 2 Oct & 0* & 0 & Gate \\
A-112 & 1.1.2 & EVM control stand-up & 7 & A-111 & 3 Oct & 9 Oct & 25 Nov & 1 Dec & 53 & 53 & N \\
A-113 & 1.1.3 & ASA / CASA / GBRMPA permits & 76 & --- & 10 Aug & 24 Oct & 10 Aug & 24 Oct & 0 & 0 & Y \\
A-114 & 1.1.4 & Juru monitor mobilisation & 62 & --- & 10 Aug & 10 Oct & 15 Aug & 15 Oct & 5 & 5 & N \\
A-121 & 1.2.1a & Site mobilisation and cleanroom & 19 & --- & 10 Aug & 28 Aug & 13 Aug & 31 Aug & 3 & 0 & N \\
A-122 & 1.2.1b & Receiving and unpacking & 8 & A-100, A-121 & 14 Sep & 21 Sep & 14 Sep & 21 Sep & 0 & 0 & Y \\
A-123 & 1.2.2 & Crane and vertical stacking & 24 & A-122 & 22 Sep & 15 Oct & 22 Sep & 15 Oct & 0 & 0 & Y \\
A-124 & 1.2.3 & Cryo GSE piping and checkouts & 45 & A-121 & 29 Aug & 12 Oct & 1 Sep & 15 Oct & 3 & 3 & N \\
A-125 & 1.2.4 & RF telemetry calibration & 18 & A-121 & 29 Aug & 15 Sep & 7 Oct & 24 Oct & 39 & 39 & N \\
A-131 & 1.3.1 & Avionics power-on sweeps & 23 & A-122 & 22 Sep & 14 Oct & 23 Sep & 15 Oct & 1 & 1 & N \\
A-132 & 1.3.2 & Wet Dress Rehearsal (cold) & 9 & A-114, A-123, A-124, A-131 & 16 Oct & 24 Oct & 16 Oct & 24 Oct & 0 & 0 & Y \\
A-133 & 1.3.3 & 5-second static fire & 6 & A-132, A-113, A-125 & 25 Oct & 30 Oct & 25 Oct & 30 Oct & 0 & 0 & Y \\
A-134 & 1.3.4 & Post-fire / pre-kitted spare & 6 & A-133 & 31 Oct & 5 Nov & 31 Oct & 5 Nov & 0 & 0 & Y \\
A-141 & 1.4.1 & LRR and RSO clearance & 5 & A-134 & 6 Nov & 10 Nov & 6 Nov & 10 Nov & 0 & 0 & Y \\
A-142 & 1.4.2 & Range exclusion and countdown & 4 & A-141 & 11 Nov & 14 Nov & 11 Nov & 14 Nov & 0 & 0 & Y \\
A-143 & 1.4.3 & Launch execution ($T$-0) & 1 & A-142 & 15 Nov & 15 Nov & 15 Nov & 15 Nov & 0 & 0 & Y \\
A-144 & 1.4.4 & Telemetry decrypt and handover & 5 & A-143 & 16 Nov & 20 Nov & 16 Nov & 20 Nov & 0 & 0 & Y \\
A-145 & 1.4.5 & Demobilisation and closeout & 11 & A-144 & 21 Nov & 1 Dec & 21 Nov & 1 Dec & 0 & 0 & Y \\
\end{longtable}
\end{landscape}

\subsection{Schedule baseline, Gantt visualisation and milestone integrity}
\label{sec:schedule-baseline}

The dates in Table~\ref{tab:precedence} are the schedule baseline. Figure~\ref{fig:gantt_chart} plots the same $ES$/$EF$ values. August and September carry permitting, heritage, mobilisation and GSE. October is the physical peak: stacking 15~October, WDR 16--24~October, static fire 25--30~October. November is the flight-clearance tail.

Two Week~4 collisions are closed here. Objective~1 dated pad integration 15~October and M-4 stacking 28~October. M-4 is now stacking on \textbf{15~October~2026}, which restores the cold-flow window before hot fire. ``Within five days of lift-off'' was also dated M-8 at 22~November, seven days after a 15~November $T$-0. M-8 is now \textbf{20~November}. The facility licence is already held, so permits gate static fire, not stacking.

Three gates sit on the baseline.
\begin{itemize}[leftmargin=1.4em,itemsep=0.25em]
\item Gate~1 is the PEP lock on 2~October.
\item Gate~2 is permits on 24~October. Without those instruments static fire cannot start. That path is already $TF = 0$.
\item Gate~3 is LRR on 10~November.
\end{itemize}
Crashing stacking without crashing the permit path therefore cannot pull $T$-0.

\begin{figure}[H]
\centering
\includegraphics[width=\linewidth]{gantt_chart.pdf}
\caption[Baseline Gantt chart]{Baseline Gantt from 10~August to 1~December~2026, plotted from the same $ES$/$EF$ values as Table~\ref{tab:precedence}. Crimson bars are launch-critical. The hatched grey bar is the external \$0 freight wait. Teal whiskers show total float to $LF$. Vertical markers are M-2 (navy, Gate~1) and M-4 through M-9. Option~C is not drawn on this baseline.}
\label{fig:gantt_chart}
\end{figure}

\subsection{Schedule compression: crashing and fast-tracking}
\label{sec:schedule-compression}

Fast-tracking overlaps sequential work, usually at zero direct cost. Crashing adds resource on critical work and is paid as a cost slope. The Board question is whether $T$-0 can move from 15~November to 10~November if R-14 is realised. Five days must come off a $TF = 0$ path that sets lift-off. After the static-fire merge that path is diagnostics through lift-off. Crashing stacking by five days saves \textbf{zero} days on $T$-0 while permits still finish on 24~October. Overlapping stacking with cryogenic GSE also saves zero days: GSE already runs in parallel.

The crash-cost slope on a crashed activity is
\begin{equation}
\text{Cost slope} = \frac{\Delta C}{\Delta T}
= \frac{\text{crash cost} - \text{normal cost}}{\text{normal duration} - \text{crash duration}}.
\label{eq:slope}
\end{equation}
Table~\ref{tab:compression} and Figure~\ref{fig:compression_curve} compare three options that each claim five days. Only Option~C is both CPM-valid and acceptable against the LRR hold.

\begin{table}[H]
\centering
\small
\setlength{\tabcolsep}{3.5pt}
\caption[Schedule compression evaluation]{Schedule compression evaluation. All three options save five days on $T$-0 if they work as drawn. Baseline $T$-0 stays 15~November. Option~C is a draw on the \$78{,}890 reserve if R-14 fires, not an addition to the \$1{,}500{,}000 base.}
\label{tab:compression}
\begin{tabularx}{\textwidth}{@{}l >{\raggedright\arraybackslash}X c r r >{\raggedright\arraybackslash}X l@{}}
\toprule
Option & Activities modified & Days & Added cost & Slope & Secondary risks & Verdict \\
\midrule
A Pure fast-track & Diagnostics $\to$ LRR SS+1 (review overlaps diagnostics) & 5 & \$0 & \$0/d & Review before diagnostic close-out; LRU rework; AS9100D hold-point failure & Rejected \\
B Pure crash & Diagnostics crashed 6~$\to$~1~d & 5 & \$40{,}000 & \$8{,}000/d & Fatigue; skipped thermal/electrical screens; recycle into LRR & Rejected \\
C Hybrid (if R-14 fires) & Diagnostics 6~$\to$~1~d extra labour (\$21{,}000); LRR kept at 5~d; countdown kept at 4~d & 5 & \textbf{\$21{,}000} & \textbf{\$4{,}200/d} & Screens kept; R-15 4~d marine hold intact & Contingent \\
\bottomrule
\end{tabularx}
\end{table}

Decision: baseline $T$-0 stays 15~November. If R-14 fires, the Board draws Option~C and lift-off moves to 10~November. That is the only option valid on the network and acceptable on quality.

Option~A is a pure fast-track. It buys five days at \$0 by starting LRR one day after diagnostics start, while those diagnostics are still open. HP-4 is a stop: flight clearance cannot start on incomplete post-fire data \citep{sae2016}. Signing early is a hold-point failure. Waiting until diagnostics finish does not save the five days; it creates rework. Option~A is rejected.

Option~B is a pure crash of diagnostics from six days to one. The slope is $\$40{,}000/5 = \$8{,}000$ per day. A 24-hour war room on a vehicle that has just fired is the work most likely to skip thermal and electrical screens. Option~B is rejected.

Option~C takes the five days only from diagnostics, and extra labour keeps the screens Option~B would skip: two specialists and two technicians for five days ($2\times\$165 + 2\times\$95 = \$520$/h $\times$ 8~h $=$ \$4{,}160/d, published at the \$4{,}200/d crawler rate). The slope is $\$21{,}000/5 = \$4{,}200$ per day. LRR stays five days and countdown stays four, so the R-15 marine hold is kept. If invoked: diagnostics 31~October, review 1--5~November, countdown 6--9~November, lift-off \textbf{10~November}. The \$21{,}000 comes from the \$78{,}890 contingency only if R-14 is realised, not from the \$1{,}500{,}000 base. Feeder float stays at 1~day (avionics) and 3~days (cryogenic GSE).

Option~C is chosen because it is cheaper than Option~B (\$4{,}200/d, not \$8{,}000/d), it keeps the LRR hold Option~A would break, and crashing stacking cannot move $T$-0 while permits finish on 24~October.

\begin{figure}[H]
\centering
\includegraphics[width=0.98\linewidth]{compression_tradeoff.pdf}
\caption[Compression cost and risk trade-off]{Compression trade-off for a five-day pull-in of $T$-0. Left: direct added cost. Right: residual schedule and quality risk after treatment. Option~C is the only CPM-valid pull-in that keeps LRR at five days, countdown at four days, and pays \$4{,}200 per day rather than \$8{,}000.}
\label{fig:compression_curve}
\end{figure}
